\documentclass[11pt,a4paper]{article}
\usepackage[T1]{fontenc}
\usepackage{lmodern,amsmath,amssymb,booktabs}
\usepackage[margin=1in]{geometry}
\setlength{\parindent}{0pt}
\setlength{\parskip}{6pt}
\pagestyle{plain}
\begin{document}
\textbf{1.}

(a)
\[
J(w_1,w_2)=\frac12(w_1-2)^2+\frac92(w_2+1)^2\geq0.
\]
Equality holds exactly at
\[
w^*=\begin{pmatrix}2\\-1\end{pmatrix},\qquad J(w^*)=0.
\]
\[
\nabla J(w^{(t)})=
\begin{pmatrix}w_1^{(t)}-2\\9(w_2^{(t)}+1)\end{pmatrix}.
\]
For $w^{(0)}=(0,0)^\top$ and $\eta=1/5$,
\[
\nabla J(w^{(0)})=\begin{pmatrix}-2\\9\end{pmatrix},\qquad
w^{(1)}=w^{(0)}-\frac15\nabla J(w^{(0)})
=\begin{pmatrix}2/5\\-9/5\end{pmatrix}.
\]

(b) The coordinate errors satisfy
\[
w_1^{(t+1)}-2=(1-\eta)(w_1^{(t)}-2),\qquad
w_2^{(t+1)}+1=(1-9\eta)(w_2^{(t)}+1).
\]
Convergence from every initial vector requires and is ensured by
\[
|1-\eta|<1,\quad |1-9\eta|<1
\quad\Longleftrightarrow\quad 0<\eta<\frac29.
\]
Starting from $(0,0)^\top$,
\[
w_1^{(t)}=2-2(1-\eta)^t,\qquad
w_2^{(t)}=-1+(1-9\eta)^t.
\]
For $\eta=1/5$,
\[
w_1^{(t)}=2-2(4/5)^t,\qquad
w_2^{(t)}=-1+(-4/5)^t,\qquad
J(w^{(t)})=\frac{13}{2}\left(\frac{16}{25}\right)^t\longrightarrow0.
\]
The first coordinate increases to $2$; the second oscillates around $-1$
with decreasing amplitude.

For $\eta=1/4$,
\[
w_1^{(t)}=2-2(3/4)^t,\qquad
w_2^{(t)}=-1+(-5/4)^t,
\]
\[
J(w^{(t)})=2\left(\frac9{16}\right)^t
+\frac92\left(\frac{25}{16}\right)^t\longrightarrow\infty.
\]
The first coordinate still converges, but the second oscillates with
increasing amplitude, so the iterates diverge.

\newpage
\textbf{2.}

(a) Let $f_k(x)=p(x\mid C_k)$. The likelihood ratio is
\[
\frac{f_2(x)}{f_1(x)}
=\exp\left(-\frac{(x-2)^2}{2}+\frac{x^2}{2}\right)
=e^{2x-2}.
\]
By Bayes' rule,
\begin{align*}
P(C_2\mid X=x)
&=\frac{(2/3)f_2(x)}{(1/3)f_1(x)+(2/3)f_2(x)}\\
&=\frac{2e^{2x-2}}{1+2e^{2x-2}}
=\frac{1}{1+e^{-(2x-2+\log2)}}.
\end{align*}
At $x=1$,
\[
P(C_2\mid X=1)=\frac23>\frac12,
\]
so the predicted class is $C_2$.

(b) Under zero-one loss, the decision boundary satisfies
\[
\frac{P(C_2\mid x)}{P(C_1\mid x)}=2e^{2x-2}=1.
\]
Thus
\[
x^*=1-\frac{\log2}{2}\approx0.653426,
\qquad
\widehat C(x)=
\begin{cases}
C_1,&x<x^*,\\
C_2,&x>x^*.
\end{cases}
\]
The posterior probabilities tie at $x=x^*$.

With equal priors,
\[
\frac{P(C_2\mid x)}{P(C_1\mid x)}=e^{2x-2},\qquad x^*=1.
\]
The boundary moves right by $(\log2)/2$: predict $C_1$ for $x<1$ and
$C_2$ for $x>1$, with a tie at $x=1$.

\newpage
\textbf{3.}

(a) There are three successes and two failures, so
\[
L(\theta)=\theta^3(1-\theta)^2,\qquad
\log L(\theta)=3\log\theta+2\log(1-\theta).
\]
\[
\frac{d}{d\theta}\log L(\theta)
=\frac3\theta-\frac2{1-\theta}=0
\quad\Longrightarrow\quad \widehat\theta_{\mathrm{ML}}=\frac35.
\]
For the prior,
\[
1=C\int_0^1\theta(1-\theta)^3\,d\theta
=C\left(\frac12-1+\frac34-\frac15\right)=\frac C{20},
\qquad C=20.
\]
The posterior density satisfies
\[
p(\theta\mid D)\propto L(\theta)q(\theta)
\propto\theta^4(1-\theta)^5.
\]
Hence
\[
\frac{d}{d\theta}\log p(\theta\mid D)
=\frac4\theta-\frac5{1-\theta}=0
\quad\Longrightarrow\quad \widehat\theta_{\mathrm{MAP}}=\frac49.
\]
Both log objectives have strictly negative second derivatives on $(0,1)$,
so these are their unique maxima.

(b) The implemented objective is
\begin{align*}
J(\theta)
&=-\frac15\bigl(3\log\theta+2\log(1-\theta)\bigr)
-\log20-\log\theta-3\log(1-\theta)\\
&=-\frac85\log\theta-\frac{17}{5}\log(1-\theta)-\log20.
\end{align*}
\[
J'(\theta)=-\frac8{5\theta}+\frac{17}{5(1-\theta)}=0
\quad\Longrightarrow\quad \widehat\theta=\frac8{25}.
\]
Since
\[
J''(\theta)=\frac8{5\theta^2}+\frac{17}{5(1-\theta)^2}>0,
\]
this is the unique minimum, and $8/25\ne4/9$.

To recover the MAP estimate with the averaged likelihood, use
\[
J_{\mathrm{MAP}}(\theta)
=-\frac15\sum_{i=1}^5\log p(y_i\mid\theta)-\frac15\log q(\theta)
=-\frac15\log\bigl(L(\theta)q(\theta)\bigr).
\]
Thus the prior term must also be divided by $5$.

\newpage
\textbf{4.}

(a) At $x=\log2$,
\[
a=(\log2,-\log2,0)^\top,\qquad
\bigl(e^{a_1},e^{a_2},e^{a_3}\bigr)=(2,1/2,1).
\]
\[
(p_1,p_2,p_3)=\left(\frac47,\frac17,\frac27\right),\qquad
\widehat y=1,\qquad \ell(a,3)=-\log\frac27=\log\frac72.
\]
For arbitrary scores and label $y$,
\[
\ell(a,y)=-a_y+\log\left(\sum_{k=1}^3e^{a_k}\right).
\]
Therefore
\[
\frac{\partial\ell}{\partial a_j}
=-\mathrm{I}[j=y]+\frac{e^{a_j}}{\sum_{k=1}^3e^{a_k}}
=p_j-\mathrm{I}[j=y].
\]

(b) For any constant $c$,
\[
p_k(a+c\mathbf1)
=\frac{e^{a_k+c}}{\sum_j e^{a_j+c}}
=\frac{e^ce^{a_k}}{e^c\sum_j e^{a_j}}
=p_k(a).
\]
Consequently,
\[
\ell(a+c\mathbf1,y)=-\log p_y(a+c\mathbf1)
=-\log p_y(a)=\ell(a,y).
\]
For $a=(1000,1001,999)^\top$, subtract $\max_j a_j=1001$:
\[
b=a-1001\mathbf1=(-1,0,-2)^\top,\qquad
D=e^{-1}+1+e^{-2}.
\]
\[
(p_1,p_2,p_3)=\frac{(e^{-1},1,e^{-2})}{D}
\approx(0.244728,0.665241,0.090031).
\]
\[
\ell(a,3)=-\log\frac{e^{-2}}D
=2+\log D\approx2.407606.
\]
Only nonpositive scores are exponentiated.

\newpage
\textbf{5.}

(a) For either class $y\in\{0,1\}$,
\[
P(X_1=0\mid y)=P(X_1=1\mid y)
=P(X_2=0\mid y)=P(X_2=1\mid y)=\frac12.
\]
For every input $(x_1,x_2)$, the Naive Bayes scores are
\[
P(Y=y)P(X_1=x_1\mid y)P(X_2=x_2\mid y)
=\frac12\cdot\frac12\cdot\frac12=\frac18.
\]
Thus both estimated posteriors equal $1/2$. The tie rule gives
$\widehat Y_{\mathrm{NB}}=0$ at all four inputs, and
\[
P(\widehat Y_{\mathrm{NB}}\ne Y)=P(Y=1)=\frac12.
\]
Using the joint distributions instead,
\[
P(Y=1\mid x_1,x_2)
=\frac{P(x_1,x_2\mid Y=1)}
{P(x_1,x_2\mid Y=0)+P(x_1,x_2\mid Y=1)}.
\]
\[
\begin{array}{c|c|c|c}
\toprule
(x_1,x_2)&P(Y=1\mid x_1,x_2)&\widehat Y_{\mathrm{NB}}&\widehat Y_{\mathrm{joint}}\\
\midrule
(0,0)&1/5&0&0\\
(0,1)&4/5&0&1\\
(1,0)&4/5&0&1\\
(1,1)&1/5&0&0\\
\bottomrule
\end{array}
\]
The joint classifier has error
\begin{align*}
P(\widehat Y_{\mathrm{joint}}\ne Y)
&=\frac12P(X_1\ne X_2\mid Y=0)
+\frac12P(X_1=X_2\mid Y=1)\\
&=\frac12(0.1+0.1)+\frac12(0.1+0.1)=\frac15.
\end{align*}

(b) For $Z=\mathrm{I}[X_1\ne X_2]$,
\[
P(Z=1\mid Y=0)=\frac15,\quad P(Z=1\mid Y=1)=\frac45,
\qquad P(Z=0)=P(Z=1)=\frac12.
\]
Hence
\[
P(Y=1\mid Z=z)=
\begin{cases}
1/5,&z=0,\\
4/5,&z=1.
\end{cases}
\]
The classifier predicts $\widehat Y=Z$ and has error
\[
P(\widehat Y\ne Y)=\frac12\cdot\frac15+\frac12\cdot\frac15=\frac15.
\]
More data cannot improve the original classifier with the stipulated exact
marginals: both classes already have identical single-feature distributions.
The discriminative information lies in the dependence between $X_1$ and $X_2$,
which the conditional-independence model discards. The feature $Z$ retains it.
\end{document}
